\textbf{COLUNI 2023 -- Questão 35}\par\medskip A conquista de maiores espaços pelas mulheres na sociedade contemporânea é reflexo de uma série de lutas e movimentos que buscaram reconhecimento e plena participação das mulheres no meio social e nas instâncias de representação política. Nas últimas décadas, tivemos várias personagens que exemplificam o movimento de ocupação feminina em cargos de liderança política nacionais. Como exemplo, temos a eleição de Dilma Rousseff em 2010, primeira mulher na presidência da República brasileira; e outros nomes importantes como Margareth Tatcher, Primeira Ministra na Inglaterra; Angela Merkel, Chanceler da Alemanha; Cristina Kirchner, Presidenta da Argentina; Golda Meir, Primeira Ministra de Israel, e Benazir Butho, Primeira Ministra do Paquistão.

\par

No Brasil, segundo dados “do IBGE, mais da metade da população brasileira (51,13\%) é feminina, e elas representam, segundo o Tribunal Superior Eleitoral, 53\% do eleitorado. No entanto, ocupam hoje menos de 15\% dos cargos eletivos”. Disponível em: https://www12.senado.leg.br/noticias/infomaterias/2022/05/aliados -na-luta-por-mais-mulheres-na- politica. Acesso em: 18 Ago. 2022.

\par

Leia as afirmativas sobre a condição feminina na sociedade contemporânea.

\par

I.   A Revolução Industrial, nos séculos XVIII e XIX, que causou significativas transformações nas relações de trabalho e no sistema de produção, propiciou às mulheres, simultaneamente, o direito ao trabalho e ao voto, colocando-as em igualdade de condições com os homens.

II. A Revolução Sexual das décadas de 1960-1970, como movimento que desafiou os códigos tradicionais de comportamento relacionados à sexualidade humana e aos relacionamentos interpessoais, estimulou a fixação da mulher no lar, proporcionando mais tempo para os cuidados com a família e com o casamento.

III. Embora ocupem cargos de liderança política, em vários países, as mulheres ainda constituem minoria entre os grupos dirigentes políticos e econômicos em diferentes espaços de poder no Brasil e no mundo.

IV. Existe até hoje uma divisão sexual do trabalho, que pode ser verificada nos mais distintos âmbitos da sociedade, gerando uma desigual repartição entre os sexos. No mercado de trabalho, de forma geral, mulheres ainda ganham salários menores que os homens, exercendo as mesmas funções e os mesmos cargos.

\par

Estão CORRETAS apenas as afirmativas:\begin{enumerate}
\item III e IV.
\item I e III.
\item II e IV.
\item I, II e IV.\end{enumerate}
